% Drop-in LaTeX corrections for the TCSSOP manuscript.
% These blocks align the stage-state and sample-feasibility description with
% the Java scheduler implementation.

% ---------------------------------------------------------------------------
% 1) Replace or expand Section 3.4, Stage-state logic and sample feasibility
% ---------------------------------------------------------------------------

The scheduling logic is implemented through four job categories: gas amount
determination, pull-down, other downstream tests, and consumer-usage tests. Gas
amount determination is executed at most once for each product when it is
required by the product's test matrix. If it is required, it is the first
operation of the product and is released at time zero. If it is not required,
the product is treated as already released for the pull-down/downstream logic.

Pull-down is executed once for each physical sample of a product. For product
\(p\) with \(s_p\) samples, the pull-down phase therefore generates \(s_p\)
sample-specific operations. Each pull-down operation can start only after the
gas amount determination operation of the same product has finished. Different
pull-down operations of the same product may run in parallel on different
samples and compatible stations, subject to station availability and
sample-level non-overlap.

The downstream non-consumer tests consist of the required energy-efficiency,
performance, freezing-capacity, and temperature-rise tests. In the implemented
configuration, these tests are released only after the full pull-down phase of
the product has been completed, that is, after all sample-level pull-down
operations have finished. Each required downstream non-consumer test is
scheduled once per product and may be assigned to any available sample of that
product after this release time.

Consumer-usage tests are modeled as the final category. A required
consumer-usage test is released only after the pull-down phase has been
completed and after all required downstream non-consumer tests of the same
product have been started. Thus, the release time for consumer-usage tests is
the maximum of the pull-down phase completion time and the latest start time
among the required downstream non-consumer tests. If a product has no required
downstream non-consumer tests, the consumer-usage release is determined by the
pull-down phase completion time, or by the gas completion time when pull-down
is not required.

This stage logic can be summarized as follows. Let \(G_p\) be the gas operation
of product \(p\), \(PD_{pj}\) be the pull-down operation of sample \(j\), and
\(\mathcal{O}_p\) and \(\mathcal{U}_p\) be the sets of required downstream
non-consumer and consumer-usage tests, respectively. If gas and pull-down are
required, then
\[
G_p \prec PD_{pj} \qquad \forall j=1,\ldots,s_p.
\]
The implemented downstream release rule is
\[
\max_{j=1,\ldots,s_p} C_{PD_{pj}} \le S_o
\qquad \forall o\in\mathcal{O}_p,
\]
where \(S_o\) denotes the start time of downstream non-consumer test \(o\). For
each consumer-usage test \(u\in\mathcal{U}_p\),
\[
S_u \ge
\max\left\{
\max_{j=1,\ldots,s_p} C_{PD_{pj}},
\max_{o\in\mathcal{O}_p} S_o
\right\}.
\]
The second term is based on the start times of the downstream non-consumer
tests, not their completion times. This reflects the implemented stage-state
rule that consumer-usage testing can be released once all relevant downstream
tests have been initiated and the pull-down phase has ended.

Sample feasibility is enforced independently from product-stage feasibility.
At any time, a physical sample can process at most one operation. Therefore,
even when multiple tests of a product are released simultaneously, they can
run in parallel only if they are assigned to different samples and compatible
stations. The scheduler updates sample availability after every operation and
uses this availability when selecting the earliest feasible station/sample
assignment.

% ---------------------------------------------------------------------------
% 2) Replace vague references to "relevant downstream-stage conditions"
% ---------------------------------------------------------------------------

% Suggested replacement sentence:
Consumer-usage tests are released after the product's pull-down phase is
complete and after all required downstream non-consumer tests have started;
they do not wait for those downstream non-consumer tests to finish.
